\chapter{Regulatory Capital for Trading Books}\label{ch:rc:regulatory-capital-for-trading-books}

Under the market-risk rules that the Basel Committee finalised in January 2019, a bank that
wants to set a desk's capital with its own model must show, every quarter, that the model
explains the desk's own P\&L. If the risk model's P\&L and the P\&L of the pricing systems
diverge too much, the desk pays a surcharge; if they diverge further, the desk loses the
model and falls back on a standardised formula that, for the book of this chapter, is almost
three times as high. The rules took years to reach the statute books: deferred once by the
Committee, postponed for two years in the European Union, delayed in the United Kingdom, and still a
proposal in the United States in 2026. This chapter computes both approaches on chapter~21's
book: the standardised sensitivities, the \omterm{def:rc:market-risk-measures:es}{expected shortfall} with its stressed calibration,
and the attribution test that decides between them.

\section{The trading book and its boundary}

\begin{definition}[Fundamental Review of the Trading Book]\label{def:rc:regulatory-capital-for-trading-books:frtb}
The \emph{Fundamental Review of the Trading Book}\index{Fundamental Review of the Trading Book}
(FRTB) is the Basel Committee's market-risk capital framework of January 2019: a revised
boundary between the books, a standardised approach built on sensitivities, and an \omterm{def:rc:regulatory-capital-for-trading-books:ima}{internal
models approach} built on \omterm{def:rc:market-risk-measures:es}{expected shortfall} with desk-level eligibility tests.
\end{definition}

\begin{definition}[Trading book, banking book]\label{def:rc:regulatory-capital-for-trading-books:books}
The \emph{trading book}\index{trading book} holds instruments held for short-term resale, for
profiting from short-term price movements, for locking in arbitrage profits, or for hedging
such positions, fair valued daily through P\&L; everything else is in the \emph{banking
book}\index{banking book}, whose market risk (mainly interest rates) is managed as in
chapter~24.
\end{definition}

The boundary matters because the two books carry different capital; moving instruments
between them to reduce capital is restricted and, where allowed, any capital benefit is
disallowed.

\section{The standardised approach: sensitivities, curvature, default and residual risks}

\begin{definition}[Sensitivities-based method]\label{def:rc:regulatory-capital-for-trading-books:sbm}
The \emph{sensitivities-based method}\index{sensitivities-based method} computes prescribed
sensitivities (a PV01 per tenor for interest rates, the value change for a 1\% move for FX),
multiplies them by prescribed risk weights, and aggregates the weighted sensitivities within
buckets ($K_b = \sqrt{\max(0,\sum_{kl}\rho_{kl}WS_kWS_l)}$) and across buckets
($\sqrt{\sum_bK_b^2+\sum_{b\ne c}\gamma_{bc}S_bS_c}$), under three correlation scenarios (the prescribed
correlations, 1.25 times them capped at one, and $\max(2\rho-1,0.75\rho)$); the requirement is the
largest of the three totals.
\end{definition}

\begin{definition}[Curvature risk charge]\label{def:rc:regulatory-capital-for-trading-books:curv}
The \emph{curvature risk charge}\index{curvature risk charge} captures the losses of options
beyond their delta: each \omterm{def:rc:market-risk-measures:factor}{risk factor} is shocked up and down by its risk weight, and the loss
beyond the delta term, $\mathrm{CVR}^\pm = -\sum_i(V_i(x^\pm)-V_i(x)\mp\mathrm{RW}\,s_i)$, is charged, the larger
of the two directions.
\end{definition}

\begin{definition}[Default risk charge, residual risk add-on]\label{def:rc:regulatory-capital-for-trading-books:drc}
The \emph{default risk charge}\index{default risk charge} captures the \omterm{def:rc:reduced-form-credit:jtd}{jump-to-default risk} of
credit and equity positions that spread shocks miss. The \emph{residual risk
add-on}\index{residual risk add-on} charges a flat share of notional for risks the
sensitivities do not capture: 1\% for instruments with exotic underlyings, 0.1\% for other
residual risks.
\end{definition}

\omcode{code/firm/frtb/firm_frtb.py}{37}{63}{Within- and across-bucket aggregation, and the GIRR delta charge under a correlation scenario.}\label{lst:rc:regulatory-capital-for-trading-books:sbm}

\begin{example}[The standardised charge of the book]\label{ex:rc:regulatory-capital-for-trading-books:sa}
Chapter~21's book, on 23 September 2026. Its GIRR sensitivities (mapping the two bonds'
par-yield sensitivities to the two- and ten-year USD tenors, a simplification) are $+56\,525$
and $-155\,012$ dollars per basis point; with risk weights $1.3\%/\sqrt2$ and $1.1\%/\sqrt2$ and a
correlation of 88.7\%, the GIRR delta charge is USD~7.83 million. FX delta (EUR and JPY, risk weight
$15\%/\sqrt2$, cross-currency correlation 60\%) is USD~5.30 million; FX vega, the straddle's vega times
its 8\% volatility at a 100\% risk weight, USD~14.49 million; FX curvature, from shocking EURUSD by
$\pm10.6\%$, USD~39.23 million. In the low correlation scenario the total is USD~68.29 million, the
largest of the three (\cref{fig:rc:regulatory-capital-for-trading-books:sa}).
\end{example}

\begin{omfigure}
\begin{tikzpicture}
\begin{axis}[width=11.5cm,height=5.4cm,ybar,bar width=8pt,
  symbolic x coords={GIRR delta,FX delta,FX vega,FX curvature},xtick=data,
  ylabel={charge (USD million)},
  tick label style={font=\small},label style={font=\small},
  ymin=0,ymax=45,grid=major,grid style={gray!25},
  legend columns=3,legend style={font=\footnotesize,at={(0.5,-0.2)},anchor=north,draw=none,fill=none,/tikz/every even column/.append style={column sep=8pt}},
  table/col sep=comma]
\addplot[fill=omMeth!60,draw=omMeth] table[x=part,y=low]{figdata/rates-credit-risk/23-regulatory-capital-for-trading-books/sa.csv};
\addplot[fill=omDef!60,draw=omDef] table[x=part,y=medium]{figdata/rates-credit-risk/23-regulatory-capital-for-trading-books/sa.csv};
\addplot[fill=omThm!60,draw=omThm] table[x=part,y=high]{figdata/rates-credit-risk/23-regulatory-capital-for-trading-books/sa.csv};
\legend{low correlation,medium,high}
\end{axis}
\end{tikzpicture}

\omcaption{Standardised charges of chapter~21's book by component and correlation scenario.
The short straddle's curvature dominates; the delta charges move with the scenario (lower
correlations remove the offset between the two- and ten-year positions), vega and a single
\omterm{def:rc:rates-risk:factors}{curvature factor} do not. Data: the chapter's tutorial.}\label{fig:rc:regulatory-capital-for-trading-books:sa}
\end{omfigure}

\section{The internal models approach: expected shortfall with liquidity horizons}

\begin{definition}[Internal models approach, liquidity horizon]\label{def:rc:regulatory-capital-for-trading-books:ima}
The \emph{internal models approach}\index{internal models approach} sets capital from the
bank's own model of \omterm{def:rc:market-risk-measures:es}{expected shortfall} at 97.5\%, computed on a ten-day base horizon and scaled
up for \omterm{def:rc:market-risk-measures:factor}{risk factors} with longer \emph{liquidity horizons}\index{liquidity horizon} (10, 20, 40,
60 or 120 days, the time assumed to exit or hedge a position without moving prices):
$\mathrm{ES} = \sqrt{\mathrm{ES}_T(P)^2+\sum_{j\ge2}\bigl(\mathrm{ES}_T(P,j)\sqrt{(\mathrm{LH}_j-\mathrm{LH}_{j-1})/T}\bigr)^2}$.
\end{definition}

\begin{definition}[Stressed expected shortfall]\label{def:rc:regulatory-capital-for-trading-books:ses}
\emph{Stressed expected shortfall}\index{stressed expected shortfall} calibrates the ES to the
most severe twelve-month period for the bank's portfolio, computed on a reduced set of \omterm{def:rc:market-risk-measures:factor}{risk
factors} with a long history and scaled up by the ratio of the current ES with all factors to
the current ES with the reduced set (at least one).
\end{definition}

The idea is older than the FRTB: the Basel Committee's July 2009 revisions (Basel~2.5) added to the VaR
charge a stressed VaR, a ten-day 99\% VaR of the current portfolio calibrated to a continuous twelve-month
period of significant stress, such as 2007--08, each charge taken as the higher of its latest value and a
multiple of its sixty-day average.

Capital for modellable factors is $\mathrm{IMCC} = \rho\,\mathrm{ES}(\text{all})+(1-\rho)\sum_i\mathrm{ES}(\text{class }i)$ with $\rho = 0.5$,
half the diversified ES and half the sum over risk classes, and the requirement is the larger of
yesterday's IMCC and 1.5 times its sixty-day average, plus the stressed charge for
non-modellable factors.

\begin{example}[The internal-model charge]\label{ex:rc:regulatory-capital-for-trading-books:ima}
On overlapping ten-day moves, the book's 97.5\% ES over the last year is USD~5.27 million. Scanning
every twelve-month window since 2016, the most severe ends on 3 October 2022
(\cref{fig:rc:regulatory-capital-for-trading-books:windows}): its ES is USD~14.41 million for the whole
book, 5.67 million for rates alone and 14.01 million for FX alone. All the factors are major-currency
rates and FX pairs with a ten-day \omterm{def:rc:regulatory-capital-for-trading-books:ima}{liquidity horizon}, so no scaling applies. The IMCC is
$0.5\times14.41+0.5\times(5.67+14.01) = \text{USD}~17.04$ million and, with the multiplier of 1.5, the charge
is USD~25.56 million, 37\% of the standardised one.
\end{example}

\begin{omfigure}
\begin{tikzpicture}
\begin{axis}[width=11.5cm,height=5.2cm,
  xlabel={end of the twelve-month window},ylabel={ten-day 97.5\% ES (USD million)},
  tick label style={font=\small},label style={font=\small},
  xmin=2017,xmax=2026.8,ymin=0,ymax=16,grid=major,grid style={gray!25},
  x tick label style={/pgf/number format/1000 sep={}},
  table/col sep=comma]
\addplot[omDef,very thick] table[x=end,y=es]{figdata/rates-credit-risk/23-regulatory-capital-for-trading-books/es_windows.csv};
\end{axis}
\end{tikzpicture}

\omcaption{\omterm{def:rc:market-risk-measures:es}{Expected shortfall} of today's book over every twelve-month window of history since
2016, by the window's end date. The stressed calibration takes the maximum, the year to October
2022. Data: US Treasury par yields, ECB reference rates; the chapter's tutorial.}\label{fig:rc:regulatory-capital-for-trading-books:windows}
\end{omfigure}

\section{Backtesting and the P\&L attribution test}

\begin{definition}[Hypothetical and risk-theoretical P\&L]\label{def:rc:regulatory-capital-for-trading-books:hpl}
The \emph{hypothetical P\&L}\index{hypothetical P\&L} (HPL) of a desk is the daily change in the
value of yesterday's positions computed by the bank's pricing systems with today's market data.
The \emph{risk-theoretical P\&L}\index{risk-theoretical P\&L} (RTPL) is the same change computed by
the risk model, from its own \omterm{def:rc:market-risk-measures:factor}{risk factors} and its own revaluation.
\end{definition}

\begin{definition}[P\&L attribution test]\label{def:rc:regulatory-capital-for-trading-books:pla}
The \emph{P\&L attribution test}\index{P\&L attribution test} compares a desk's HPL and RTPL over
250 days by the Spearman correlation of their ranks and the Kolmogorov--Smirnov distance between
their distributions: green if the correlation exceeds 0.80 and the distance is below 0.09, red if
the correlation is below 0.70 or the distance above 0.12, amber otherwise. Red desks use the
standardised approach; amber desks pay a surcharge of half the gap between the standardised and
internal-model charges (for a bank whose only desk is amber).
\end{definition}

\omcode{code/firm/frtb/firm_frtb.py}{136}{159}{Spearman's rank correlation, the Kolmogorov--Smirnov distance and the attribution zones.}\label{lst:rc:regulatory-capital-for-trading-books:pla}

\begin{example}[Three risk models of the same desk]\label{ex:rc:regulatory-capital-for-trading-books:pla}
Over the last 250 days, a risk model that revalues the book with the delta--gamma expansion
passes easily (Spearman 1.00, KS 0.008). One that has no yen factor passes too (0.956 and 0.060).
One that has no two-year factor is amber through its distribution (0.949 and 0.100). One that
proxies the ten-year yield by the two-year, a common response to a factor judged non-modellable,
is amber through its correlation: 0.703 and 0.044 (\cref{fig:rc:regulatory-capital-for-trading-books:pla}).
\end{example}

\begin{omfigure}
\begin{tikzpicture}
\begin{axis}[width=8cm,height=6.5cm,
  xlabel={hypothetical P\&L (USD million)},ylabel={risk-theoretical P\&L (USD million)},
  tick label style={font=\small},label style={font=\small},
  xmin=-2.5,xmax=2,ymin=-2.5,ymax=2,grid=major,grid style={gray!25},
  table/col sep=comma]
\addplot[omExo,dashed,domain=-2.5:2,samples=2,forget plot] {x};
\addplot[omThm,only marks,mark=*,mark size=0.9pt] table[x=hpl,y=rtpl]{figdata/rates-credit-risk/23-regulatory-capital-for-trading-books/pla.csv};
\end{axis}
\end{tikzpicture}

\omcaption{Daily \omterm{def:rc:regulatory-capital-for-trading-books:hpl}{hypothetical P\&L} of the book against the P\&L of a risk model that proxies the
ten-year yield by the two-year, over 250 days. The ranks agree only loosely (Spearman 0.703):
amber. Data: US Treasury, ECB; the chapter's tutorial.}\label{fig:rc:regulatory-capital-for-trading-books:pla}
\end{omfigure}

\section{Non-modellable risk factors}

\begin{definition}[Risk-factor eligibility test, non-modellable risk factor]\label{def:rc:regulatory-capital-for-trading-books:nmrf}
The \emph{risk-factor eligibility test}\index{risk-factor eligibility test} admits a \omterm{def:rc:market-risk-measures:factor}{risk factor}
into the internal model if it has at least 24 real price observations in the past year with no
90-day period holding fewer than four (or at least 100 in the year). A factor that fails is a
\emph{non-modellable risk factor}\index{non-modellable risk factor}: it is capitalised separately
by its own stress scenario, with little diversification.
\end{definition}

\begin{dated}{2026-09}{When the FRTB applies}\label{dat:rc:regulatory-capital-for-trading-books:dates}
The Basel Committee's January 2019 standard was to apply from 1 January 2022; in March 2020 its
oversight body deferred it to 1 January 2023. In the European Union, after the rules were
postponed by two years, they apply from 1 January 2027 with targeted, temporary adjustments for three
years, including a capital multiplier (Commission delegated act adopted on 4 June 2026, published on
1 September 2026). In the United Kingdom, the PRA consulted in June 2026 on its internal model rules,
to apply from 1 January 2028 with the attribution test's monitoring period extended from one to three
years; its other Basel~3.1 rules apply from January 2027. In the United States, the agencies proposed
in March 2026 a framework implementing the final Basel~III components, with market-risk rules for banks
with significant trading activity; comments closed on 18 June 2026.
\end{dated}

\section{Tutorial: two capital numbers for one desk}\label{tut:rc:regulatory-capital-for-trading-books:tutorial}

\begin{tutorial}
\textbf{Goal.} Compute the standardised and internal-model charges of chapter~21's book and run
the attribution test on several risk models. \textbf{End state:} the numbers of
\cref{ex:rc:regulatory-capital-for-trading-books:sa,ex:rc:regulatory-capital-for-trading-books:ima,ex:rc:regulatory-capital-for-trading-books:pla}
and the three charts.
\begin{tutsteps}
\item \textbf{Sensitivities}: \texttt{sensitivities()} bumps the book as the standard prescribes.
\item \textbf{Standardised}: \texttt{standardised()} under the three scenarios.
\item \textbf{Internal models}: \texttt{internal\_models()} with the stressed window.
\item \textbf{Attribution}: \texttt{pla(variant)}; \texttt{fig\_rc\_frtb.py} writes the charts.
\end{tutsteps}
\textbf{What to change next.} Halve the straddle and see which charge moves most; give the
yen a twenty-day \omterm{def:rc:regulatory-capital-for-trading-books:ima}{liquidity horizon} and apply the scaling formula.
\end{tutorial}

\section{Build: the FRTB engine}\label{bld:rc:regulatory-capital-for-trading-books:frtb}

\begin{build}
\bfield{Purpose} The firm's market-risk capital under both approaches and its desk eligibility
tests, fed by the risk engine of chapter~29.
\bfield{Interface} \texttt{girr\_rho}, \texttt{scenario}, \texttt{within}, \texttt{across};
\texttt{girr\_delta}, \texttt{fx\_delta}, \texttt{fx\_vega}, \texttt{cvr}, \texttt{fx\_curvature},
\texttt{sbm}; \texttt{es}, \texttt{es\_liquidity}, \texttt{imcc}, \texttt{ima\_capital},
\texttt{amber\_surcharge}; \texttt{spearman}, \texttt{ks\_stat}, \texttt{pla\_zone};
\texttt{rfet}.
\bfield{Rules} Parameters as in the Basel standard of January 2019; GIRR and FX risk classes
only; specified currencies with the $\sqrt2$ reduction; the largest scenario total.
\bfield{Acceptance tests} \texttt{code/firm/frtb/tests/}: the standard's correlation example
(88.69\%); scenario transforms; aggregation with its fallback; curvature signs; capital formula
and surcharge; liquidity scaling; attribution statistics and zone boundaries; eligibility test.
\bfield{Stretch} Credit spread, equity and commodity classes; the \omterm{def:rc:regulatory-capital-for-trading-books:drc}{default risk charge}; the
\omterm{def:rc:regulatory-capital-for-trading-books:drc}{residual risk add-on}; non-modellable stress scenarios; national variations (EU multipliers).
\end{build}

\begin{omsources}
\item Basel Committee on Banking Supervision, \emph{Minimum capital requirements for market
risk}, January 2019.
\item Bank for International Settlements, press release of 27 March 2020 on the deferral of Basel~III.
\item European Commission, ``Commission adopts temporary adjustments to Basel~III market risk
rules'', 4 June 2026.
\item Bank of England, ``PRA sets out adjustments to its market risk internal model approach
under Basel~3.1'', 19 June 2026.
\item Board of Governors of the Federal Reserve System, press release of 19 March 2026 on
proposals to modernise the regulatory capital framework.
\end{omsources}

\section{Exercises}

\begin{exercise}[$\star$]\label{exo:rc:regulatory-capital-for-trading-books:1}
Compute the GIRR correlation between the two- and ten-year tenors and its high and low scenario
values.
\end{exercise}

\begin{exercise}[$\star$]\label{exo:rc:regulatory-capital-for-trading-books:2}
Why did the low correlation scenario give the largest total for this book?
\end{exercise}

\begin{exercise}[$\star$]\label{exo:rc:regulatory-capital-for-trading-books:3}
A desk's attribution statistics are 0.82 and 0.10. What zone is it in, and what happens?
\end{exercise}

\begin{exercise}[$\star\star$]\label{exo:rc:regulatory-capital-for-trading-books:4}
Why is the internal-model charge so much lower than the standardised one for this book?
\end{exercise}

\begin{exercise}[$\star\star$]\label{exo:rc:regulatory-capital-for-trading-books:5}
A \omterm{def:rc:market-risk-measures:factor}{risk factor} has 30 observations in the year, all in its first four months. Is it modellable?
\end{exercise}

\begin{exercise}[$\star\star$]\label{exo:rc:regulatory-capital-for-trading-books:6}
Why does the IMCC mix the diversified ES with the sum of risk-class ES?
\end{exercise}

\begin{exercise}[$\star\star\star$]\label{exo:rc:regulatory-capital-for-trading-books:7}
\emph{Coding.} Compute the ES of a portfolio whose \omterm{def:rc:market-risk-measures:factor}{risk factors} split into a ten-day group with
ES 5 and a twenty-day group with ES 3 (where the second is also counted in the first), using the
liquidity-horizon formula.
\end{exercise}

\begin{exercise}[$\star\star\star$]\label{exo:rc:regulatory-capital-for-trading-books:8}
\emph{Find the flaw.} ``Our internal model gives a third of the standardised capital; the model is
better, so the regulator should let us use it for every desk.''
\end{exercise}

\section{Problem: The Desk That Failed Attribution}

\begin{problem}[{Weekend problem --- amber}]\label{pb:rc:regulatory-capital-for-trading-books:1}
Chapter~21's book is the only desk of a bank using the \omterm{def:rc:regulatory-capital-for-trading-books:ima}{internal models approach}. Its risk model
has no reliable ten-year data and proxies the ten-year yield by the two-year.

\medskip\noindent\textbf{Part I --- Two numbers.}
\begin{enumerate}
\item Give the standardised charge and the scenario that sets it.
\item Give the internal-model charge and its ingredients.
\item Which component drives the standardised charge?
\item Why is the stressed window the year to October 2022 for this book?
\item What would a longer FX volatility \omterm{def:rc:regulatory-capital-for-trading-books:ima}{liquidity horizon} change?
\end{enumerate}

\medskip\noindent\textbf{Part II --- The test.}
\begin{enumerate}[resume]
\item Give the Spearman correlation and KS distance of the proxy model.
\item Give the zone and why.
\item Give the surcharge and the resulting capital.
\item What would red mean for the desk's capital?
\item Which of the other risk models would pass, and which fail?
\end{enumerate}

\medskip\noindent\textbf{Part III --- Fixing it.}
\begin{enumerate}[resume]
\item Why does the proxy fail the rank correlation but not the distribution test?
\item What data would make the ten-year yield modellable?
\item Would a better proxy (a regression on the two-year) help?
\item What does a desk gain from a longer monitoring period, as the PRA proposes?
\item How should the desk's business decide between the fix and the surcharge?
\end{enumerate}

\medskip\noindent\textbf{Part IV --- Judgement.}
\begin{enumerate}[resume]
\item Why test attribution at all when backtesting exists?
\item What incentives does the amber surcharge create?
\item How do the EU multipliers and the UK delay change a global bank's choices?
\item State the \emph{named result}: the desk's internal-model and standardised capital, and the
Spearman and KS values that put it in amber.
\item In one sentence: what does the attribution test really check?
\end{enumerate}
\end{problem}

\section{Interview questions}

\begin{interviewq}[$\star$ \iqroles{risk, bank}]\label{iq:rc:regulatory-capital-for-trading-books:1}
What changed from Basel~2.5 to the FRTB?
\end{interviewq}

\begin{interviewq}[$\star\star$ \iqroles{risk, developer}]\label{iq:rc:regulatory-capital-for-trading-books:2}
Walk through the \omterm{def:rc:regulatory-capital-for-trading-books:sbm}{sensitivities-based method} for an interest rate swap book.
\end{interviewq}

\begin{interviewq}[$\star\star$ \iqroles{risk}]\label{iq:rc:regulatory-capital-for-trading-books:3}
What is the \omterm{def:rc:regulatory-capital-for-trading-books:pla}{P\&L attribution test} and why do desks fail it?
\end{interviewq}

\begin{interviewq}[$\star\star$ \iqroles{researcher, risk}]\label{iq:rc:regulatory-capital-for-trading-books:4}
Why \omterm{def:rc:market-risk-measures:es}{expected shortfall} with \omterm{def:rc:regulatory-capital-for-trading-books:ima}{liquidity horizons} rather than a ten-day VaR?
\end{interviewq}

\begin{interviewq}[$\star\star\star$ \iqroles{bank}]\label{iq:rc:regulatory-capital-for-trading-books:5}
How would you decide which desks should apply for internal-model approval?
\end{interviewq}

\begin{interviewq}[$\star\star\star$ \iqroles{developer}]\label{iq:rc:regulatory-capital-for-trading-books:6}
What data and systems does a bank need to run the \omterm{def:rc:regulatory-capital-for-trading-books:ima}{internal models approach}?
\end{interviewq}
